% 阴影面积的意义——物理量的累积效应
% keys 速度、电荷量、平均物理量
% license Usr
% type Tutor

\begin{issues}
\issueTODO 缺图像
\end{issues}

\subsection{物理量的累积效应}

\subsubsection{速度}
在初中里，我们已经学过“速度”，也就是单位时间内的路程。然而这个概念的局限很多，一方面，路程和时间都是数字，“速度”很自然的，也没有给出其他更多的信息。另一方面，当我们用初中定义去描述任何对象的快慢时，都是默认其快慢不变的。既然小明从西到东的速度是$6 \mathrm{m/s}$，那么在这个过程中，在任意的$0.02\mathrm s$内，移动的路程便总是$1.2\mathrm m$，这显然是不够贴近实际的。用“位移”代替“路程”，我们便可以解决第一个问题，这样定义的速度我们称之为平均速度。新增了方向信息后，那如何描述任意时刻的快慢呢？

回想一下，在我们用“路程-时间”图像描述物体运动时，匀速直线运动是一条斜率不变的直线，“匀速”便是路程均匀变化。受此启发，我们不妨用\textbf{斜率}来表示物体在某个时刻的速度，称之“\textbf{瞬时速度}”。这也是因为对于位移-时间图像，你总可以作出函数图像在任意一点的切线，该切线总是唯一的，所以定义切线斜率就是此时此刻的瞬时速度，显然，瞬时速度是平均速度在某个时刻附近的近似。
%瞬时速度和平均速度
当速度均匀变化时，我们称之为匀变速直线运动，并把速度-时间图像里的斜率称之为加速度。物体并非总是做匀变速直线运动。更符合实际的情况是变加速，即加速度总是变化的运动。
%匀变速和变加速的时间图像

匀速直线运动的速度时间图像是一条平行于时间轴的直线。对于任意一个时间段，时间区间与速度的乘积显然就是位移。但对于一般的速度-时间图像，这就很难成立了。从定义出发，你不知道哪个是特殊的“平均速度”，毕竟，平均速度其实是知晓位移后导出的，你怎可本末倒置地去找这个特殊值呢？但我们可以据匀速的情况猜测，是否任意一个时间段内，速度函数与两个时间轴之间的阴影面积就是这一段区间的位移。
\begin{figure}[ht]
\centering
\includegraphics[width=6cm]{./figures/25c441ef3a6dd49c.png}
\caption{用矩形近似阴影面积} \label{fig_area_1}
\end{figure}

\subsubsection{冲量}
\subsubsection{电荷量}
\subsection{“平均物理量”}